\scanpage{88}
% archive foliation: 90
\[
X(I)/V(I)\simeq X(J)\qquad\text{où }J=J(I)\text{ comme dessus}
\]
Ceci dit soit \deleted{$I\hookrightarrow X$}\inserted{une opération de $V$ sur $X$}
\deleted{on a donc \unsure{de} $X$ une opération $V=V(I)$}\inserted{$V$}$\hookrightarrow G_X$, d'où un $I(V)$ $V$-torseur
des bifaces de $\operatorname{Oc}(V)$, et on a \deleted{$X\simeq$}
\[
X\simeq X(I(V))\qquad\text{iso. canonique},
\]
(et non $X\simeq X(I)$). Mais pour \deleted{un} deux torseurs
$T,T'$ sous un $=$ $V$, on aura
\[
X(T\wedge T')=X(T)\wedge_V T'
\]
\margin{\figure{\begin{tikzpicture}\draw (0,0)--(0,.65);\end{tikzpicture}} à expliciter}
Soit donc
\[
\begin{gathered}
\uline t=\unsure{\delta|\tau}(I,X)\text{ le }V\text{-torseur tel que}\quad I(V)\simeq I\wedge_V\delta\\
\text{i.e.}\qquad\delta\simeq I(V)\wedge I^{-1}.
\end{gathered}
\]
On aura
\[
\boxed{X=X(I(V))=X(I)\wedge_V\uline t}
\]
\deleted{Donc et l'on} donc
\[
\boxed{X/V\simeq X(I)/V\simeq\unsure{Z|X}(J)}
\]

En résumé :

\figure{\begin{tikzpicture}\draw (0,0)--(0,2);\end{tikzpicture}}
\uline{Cor à Prop} \deleted{La donnée, d'un \unsure{gr.} $V\simeq D_2$}\inserted{Pour $I\in\operatorname{Ob}\operatorname{Ens}_4$}
donné, d'où $J$ et
$V=V(J)$, la donnée d'un $X$ et d'un plongement
$V\hookrightarrow G_X$ équivaut à la donnée d'un torseur $\uline t$ sous $V$,
via
\diagram{\begin{tikzcd}[ampersand replacement=\&]
\delta \arrow[r,mapsto] \& X(I)\wedge_V\deleted{\unsure{\delta}}\uline t
\end{tikzcd}}
On a donc
\[
X/V\simeq\unsure{Z|X}(J)\qquad\text{(iso. canonique)}.
\]

\uline{\unsure{Cor à Cor}}, \inserted{Donc} la donnée de $(X,I\hookrightarrow X)$ équivaut à
celle de a) un pt $\zeta$ de $Z(J)^*$, b) un $V$-torseur $\uline t$,
\deleted{c) un $V$-isom \unsure{un tel} \unsure{point} $X=X(I)\wedge_V\uline t$,}
et
\deleted{$X_\zeta=\text{fibre de}$}
\[
X_\zeta=X(I)_\zeta\wedge\uline t\simeq I\qquad\text{i.e.}\quad\uline t=I\wedge X(I)_\zeta^{-1}.
\]
